(1) 势流满足 $\phi_{xx}+\phi_{zz}=0$，底部 $\phi_z(-h)=0$。对 $e^{i(kx-\omega t)}$ 取 $\phi=A\cosh[k(z+h)]e^{i(kx-\omega t)}$，满足底部条件。表面位移 $\eta=\eta_0e^{i(kx-\omega t)}$。

(2) 表面运动条件 $\eta_t=\phi_z(0)$，动力条件 $\phi_t+g\eta-(\sigma/\rho)\eta_{xx}=0$。消去幅度得 $\boxed{\omega^2=(gk+\sigma k^3/\rho)\tanh(kh)}$。深水 $kh\gg1$ 时 $\omega^2=gk+\sigma k^3/\rho$；浅水 $kh\ll1$ 时 $\omega^2\simeq(gh)k^2+(\sigma h/\rho)k^4$。

(3) 深水相速度平方 $c_p^2=g/k+(\sigma/\rho)k$，求导 $-g/k^2+\sigma/\rho=0$，所以 $\boxed{k_*=(\rho g/\sigma)^{1/2}}$、$\boxed{c_{p,\min}=(4g\sigma/\rho)^{1/4}}$。在相速极值 $d(\omega/k)/dk=0$，严格等价于 $d\omega/dk=\omega/k$，因此 $c_g=c_p$。

(4) 若 $\sigma\to0$，深水相速度随波数单调下降、不再出现有限波数极小；浅水长波 $kh\ll1$ 时 $\omega\simeq\sqrt{gh}\,k$，相速与群速都趋 $\sqrt{gh}$。极短毛细波对流体黏性、表面污染、表面张力频散及线性小坡度假设十分敏感。
